\cotefile{74}
\scanpage{2}
\textcircled{I}
% Dark blemish near the upper edge also appears in mirror position on scan 3.
\[
\left\{\begin{aligned}
&I\ \text{ens de card }3.\\
&(E_i)_{i\in I}\quad\text{famille d'ensembles}\\
&\Gamma\subset\prod_{i\in I}E_i\\
&\text{Axiome : si }\deleted{\unsure{i,j\in I,\ i\ne j}}\quad i,j,k\in I\quad\text{distincts (donc }I=\{i,j,k\}\text{)}\\
&\text{alors }\ \forall x_i\in E_i,\ x_j\in E_j\ .\quad\exists\unsure{1|!}\ x_k\in E_k\quad\text{tel que }(x_i,x_j,x_k)\in\Gamma .
\end{aligned}\right.
\]
Catégorie $\unsure{\mathscr T|\mathscr C|\mathcal C}$ ($I$ variable\unsure{,|;} \unsure{mais|puis|pour} iso sur $I$)

\textcircled{II}\quad $I$ ens de card $3$

$(G_i)_{i\in I}$ famille de groupes

\diagram{\begin{tikzcd}[ampersand replacement=\&,row sep=small,column sep=small]
\& E_i\arrow[dl,dash,"G_k"']\arrow[dr,dash,"G_j"]\&\\
E_j\arrow[rr,dash,"\unsure{G_i|G_j}"']\&\& E_k
\end{tikzcd}}
\deleted{$\unsure{\tau_{ij}|E_{ij}}\colon L(G_i)\to L(G_j)$ syst. trans. d'isom. extérieurs} entre ces groupes

$(E_i)_{i\in I}$ famille d'ens

$\forall i\in I$, les groupes $\prod_{j\in I-\{i\}}G_j$ opèrent sur $E_i$, de façon que \inserted{si $I=\{i,j,k\}$} $E_i$ devienne un \uline{bitorseur} sous $G_j,G_k$ quand on renverse l'opération de $G_k$ en l'opération \unsure{opposée|inverse}.

Donc $\displaystyle\prod_{\unsure{i|0}\in I}\prod_{j\in I-\{i\}}G_j=\prod_{j\in I}G_j^{I-\{j\}}$ opère sur $P=\prod E_i$, donc $\prod_{j\in I}G_j$ opère diagonalement \inserted{par $\delta$}. Soit $x=(x_i)_{i\in I}\in P$, alors

\[
\begin{aligned}
&E_i\wedge E_j\wedge E_k &&x_i\wedge x_j\wedge x_k= \big\{ &&(g_j x_i,x_j,x_k g_j)\\
&G_j\quad G_k\quad G_i\quad G_j &&
g_j x_i g_k^{-1}\wedge g_k^{\unsure{+|a}}x_jg_i^{-1}\wedge g_i^{+}x_k g_j^{-1}
&&=(x_i g_k^{-1},g_k^{+}x_jg_i^{-1},g_i^{+}x_k)
\end{aligned}
\]

si $I=\{i,j,k\}$, les conditions suivantes sont équivalentes

\deleted{il revient au \unsure{=|même} de}

\[
\begin{aligned}
a_1)&\quad\delta(G_i)x\subset\unsure{\Gamma|\delta}(G_j\times G_k)\cdot x\\
a_2)&\quad\delta(G_j)x\subset\delta(G_k\times G_i)\cdot x\\
a_3)&\quad\Gamma(G_k)x\subset\Gamma(G_i\times G_j)\cdot x .
\end{aligned}
\]

\deleted{\unsure{Soit} \unsure{$x_i$}} On peut encore l'exprimer de la façon suivante : $\forall$ couple $(j,k)$ \deleted{\unsure{de|d'}} $j,k\in I$, $j\ne k$, on a

$\tau_{kj}\colon E_i\to\operatorname{Isom}(G_j,G_k)$ \quad [ l'image est un
% archive sheet number: 1
\scanpage{3}
% letterhead
UNIVERSITE DES SCIENCES\\
ET TECHNIQUES DU LANGUEDOC\\
\uline{\textit{SERVICE COMPTABILITE}}\\

Montpellier, le \unsure{-2|2} MAI 1977\\

% printed form
LE DIRECTEUR DE L'U.E.R. N° \handwritten{5}\\
à\\
M. LE PROFESSEUR \handwritten{GROTHENDiECK}\\

\uline{OBJET} : \textit{Ouverture de crédits \textquotedbl{}Recherche\textquotedbl{}.}\\

Mon cher Collègue,\\

Compte tenu des crédits actuellement notifiés par le Secrétariat d'Etat aux Universités, j'ai l'honneur de vous faire connaître qu'une première attribution de ..................\handwritten{10.450}.....F. vous est accordée pour l'année 1977. Votre code recherche est \fbox{\handwritten{M}\,\textbar{}\,\handwritten{022}} (suivi du numéro du bon de commande - 4 chiffres maximum).

Veuillez agréer, mon cher Collègue, l'assurance de mes sentiments les meilleurs.\\

LE DIRECTEUR DE L'U.E.R. N°.\handwritten{5}.....\\
\handwritten{\unsure{A. Micali|A. Mical}}\\
% archive sheet number: 2
\scanpage{4}
classe \inserted{extérieure} d'\deleted{isom.} \unsure{exactement}] et la condition est

\diagram{\begin{tikzcd}[ampersand replacement=\&,row sep=tiny,column sep=tiny]
\& \deleted{\unsure{x_i}}\arrow[dl,dash,"\deleted{G_k}"']\arrow[dr,dash,"\deleted{G_j}"]\&\\
{}\arrow[rr,dash,"\deleted{G_i}"']\&\&{}
\end{tikzcd}}
\diagram{\begin{tikzcd}[ampersand replacement=\&,row sep=small,column sep=small]
G_k\arrow[dr,"{\tau_{ik}(x_j)}"']\&\&G_j\arrow[ll,"{\tau_{kj}(x_i)}"']\\
\&G_i\arrow[ur,"{\tau_{ji}(x_k)}"']\&
\end{tikzcd}}
\[
\begin{aligned}
&\tau_{ji}(x_k)\tau_{ik}(x_j)\tau_{kj}(x_i)\\
&\hspace{4em}=\mathrm{id}_{G_j}
\end{aligned}
\]
(condition qui est bien invariante par permutations circulaires)

L'\uline{existence} d'un tel $x=(x_i)_{i\in I}$ équivaut à la condition que les isom. \uline{extérieurs} $\widetilde\tau_{ij}=\widetilde\tau_{ij}(x_k)$ (\deleted{$i,j,k\in I$} $i,j,k\in I$ distincts) définis par les $E_k$ ($k\in I$) entre les $G_i$ forment un syst. transitif d'isom de liens.

La condition envisagée sur $x$ ne dépend que de son orbite modulo $\prod G_i$, \deleted{Soit}\inserted{qui} est aussi son orbite sous l'un quelconque des sous-gr. $\prod_{j\in I-\{i\}}G_j$.

\margin{\(\begin{aligned}
g_i\gamma_i x&=(g_j,g_k)x\\
g_i(\gamma_i,\gamma_j,\gamma_k)x&=(\gamma_j,\gamma_k)g_i\gamma_i x=(\gamma_j g_j,\gamma_k g_k)x
\end{aligned}\)}
Soit $\Gamma$ une telle orbite, qui est donc un torseur sous l'un quelconque des \inserted{trois} groupes $\prod_{j\in I-\{i\}}G_j$. On a ainsi \inserted{une} \deleted{\unsure{certaine|certain|contenu}} partie
\[
\Gamma\subset\prod_{i\in I} E_i=P
\]
Celle-ci satisfait \deleted{\unsure{à}}\inserted{aux} conditions de \textcircled{I}.

\[
\begin{aligned}
a&=(a_i,a_j,a_k)&x_i&=a_i g_k^{-1}& &\text{détermine }g_k\text{ en fonction de }a\text{ et }x_i\\
x&=(x_i\ \ x_j\ \ \underset{?}{x_k})&x_j&=g_k a_j g_i^{-1}& &\text{détermine }g_i\text{ en fonction de }a\text{ et }x_i,x_j\\
&&x_k&=g_i a_k&&\text{détermine }x_k
\end{aligned}
\]
% archive sheet number: 3
\scanpage{5}
% letterhead
UNIVERSITE DES SCIENCES\\
ET\\
\uline{TECHNIQUES DU LANGUEDOC}\\

Montpellier, le 9 Juin 1976\\

\uline{Service du Personnel}\\
\uline{Diffusion} :\\
Tous enseignants de\\
l'Enseignement Supérieur\\
n° 3 718\\

MR LE PROFESSEUR GROTHENDIECK A.\\
MATHEMATIQUES\hfill 5\\

NOTE CONCERNANT LA PRIME DE RECHERCHE\\
\uline{SCIENTIFIQUE}

\textit{(Attribution individuelle au titre de l'année 1977)}

Tous les membres du personnel enseignant de l'Enseignement Supérieur sont invités à passer au Service du Personnel (Accueil) afin de remplir et signer les imprimés relatifs à la prime de recherche scientifique pour l'année 1977 :

à partir du \uline{13 Juin} et jusqu'au \uline{8 Juillet 1977} (date limite)

L'attention des intéressés est appelée sur le caractère obligatoire de cette formalité, le paiement de la prime de recherche ne pouvant être effectué que sur le vu des imprimés remplis et signés par les bénéficiaires.

\uline{Les délais mentionnés ci-dessus sont impératifs}.

Le Président,\\
Le Vice - Président,\\

B. CHARLES\\
% archive sheet number: 4
\scanpage{6}
satisfaisant
\[
u_i(1)=1\qquad i=1,2,3
\]
et
\[
\boxed{xyz=1\quad\Longrightarrow\quad u_1(x)u_2(y)u_3(z)=1}\tag*{\unsure{(b)|(6)|($*$)}}
\]

Prouvons que les $u_i$ sont \deleted{des hom de} \inserted{égaux à un seul} et unique hom de groupes.

Faisons $x=1$, on trouve que $u_3(y^{-1})=u_2(y)^{-1}$

Donc $u_3(y)=u_2(y^{-1})^{-1}$, pour la $=$ raison $u_1(y)=u_2(y^{-1})^{-1}$, donc $u_1(y)=u_3(y)$ et de $=$ $u_1(y)=u_2(y)$. Donc $u_1=u_2=u_3$, et $u(y)=u(y^{-1})^{-1}$, donc $\deleted{\unsure{y|u}}u(y^{-1})=u(y)^{-1}$.

On tire maintenant de \unsure{a)|a,|($*$)}
\[
\begin{aligned}
u((xy)^{-1})^{-1}&=u(x)u(y)\\
\rotatebox{90}{$=$}&\\
u(xy)&
\end{aligned}
\]
donc \qquad $u(xy)=u(x)u(y)$ \qquad cqfd.

\uline{Question} Le foncteur $\varphi$ est il une équivalence ?
% archive sheet number: 5
\scanpage{7}
Catégorie $\unsure{\mathscr T|\mathscr C|\mathcal C}'$ ($I$ variable) \unsure{puis|pour} iso sur $I$)

On a un foncteur
\[
\mathscr T'\xrightarrow{\varphi}\mathscr T
\]

\uline{Prop} Ce foncteur est pl. fidèle (et identifie donc \inserted{(à équiv. près)} $\mathscr T'$ à une sous-catégorie pleine de $\mathscr T$)

Comme ce sont des catégories fibrées sur le groupoïde des ens $I$ de card $3$, on est ramené à prouver que \inserted{le foncteur induit sur} les fibres sur $I=\{1,2,3\}$ \deleted{\unsure{sont équivalentes}} est \deleted{une \unsure{équivalence de}} \inserted{pl. fid.} La \uline{fidélité} est claire

Soit \inserted{$(E_i,G_i,\ldots,\Gamma)$} un objet de $\mathscr T'$, et $x$ \inserted{$=(x_1,x_2,x_3)$} un pt de $\Gamma$. Donc les $G_i$ sont isomorphes entre eux, \inserted{soit $G$,} et isomorphes aux $E_i$, et $\prod E_i$ s'identifie à $\deleted{\unsure{\prod}}G^3$, et $\Gamma$ à l'ens des $(x,y,z)$ tels que
\[
xyz=1
\]

Soit un autre objet \unsure{$\underline x$} $(E'_i,G'_i,\ldots,\Gamma')$ et $u$ un morphisme du \inserted{$\varphi$ du} premier dans le $\varphi$ du second, appliquant $\Gamma$ dans $\Gamma'$, soit $(x'_1,x'_{\unsure{2|i}},x'_3)$ l'image de $x=(x_1,x_2,x_3)=(1,1,1)$. Elle permet d'identifier le deuxième objet à $G'{}^3$, ($G'$ un groupe), $\Gamma'$ à l'ens des $(x',y',z')$ tels que
\[
x'y'z'=1 .
\]

Donc on a trois applications
\[
u_1,u_2,u_3\colon G\longrightarrow G'
\]
% archive sheet number: 6
\scanpage{9}
$\deleted{x\quad y}\qquad a_1\quad \deleted{\unsure{a}}\unsure{b|a}_2\quad a_3$

\[
\begin{aligned}
\deleted{\unsure{E}} \qquad E_1&\xrightarrow{\alpha_1}\operatorname{Isom}(E_2,E_3)&
E_2&\overset{\alpha_2}{\text{---}}\operatorname{Isom}(E_3,E_1)\\
&&E_3&\overset{\alpha_3}{\unsure{\text{---}|\longrightarrow}}\operatorname{Isom}(E_1,E_2)
\end{aligned}
\]

\diagram{\begin{tikzcd}[ampersand replacement=\&,row sep=large,column sep=small]
E_1\arrow[rr,"{\alpha_3(a_3)}"]\arrow[dr,bend left=16,dashed,"{u_2}"]\&\&
E_2\arrow[ll,bend left=13,dashed,"{u_3}"]\arrow[dl,"{\alpha_1(a_1)}"']\\
\&E_3\arrow[ul,"{\alpha_2(\unsure{a_2|a_1})}"']\arrow[ur,bend left=17,dashed,"{u_1}"]\&
\end{tikzcd}}

\uline{NB}
\[
\begin{aligned}
\alpha_3(a_3)&:a_1\ \text{---}\ a_2\\
\alpha_1(a_1)&:a_2\ \text{---}\ a_3\\
\alpha_2(a_2)&:a_3\ \text{---}\ a_1
\end{aligned}
\]
\deleted{\unsure{P}} Posons
\[
\begin{aligned}
\alpha_2(a_2)\alpha_1(a_1)&=u_3:E_2\ \text{---}\ E_1\\
\alpha_3(a_3)\alpha_2(a_2)&=u_1:E_3\ \text{---}\ E_2\\
\alpha_1(a_1)\alpha_3(a_3)&=u_2:E_1\ \text{---}\ E_3
\end{aligned}
\]

\uline{Admettons} \qquad $\boxed{u_1 u_2 u_3=\mathrm{id}}$

\margin{NB $\begin{aligned}[t]u_3(a_{\unsure{2|1}})&=a_1\\u_1(a_3)&=a_2\\u_2(a_1)&=a_3\end{aligned}$}

D'où syst transitif d'isom entre les $E_i$,
% archive sheet number: 8
\scanpage{10}
\uline{Liste des chercheurs 1977 =}

\uline{\unsure{Pr|P\textsuperscript{r}} Grothendieck.}

\diagram{$\begin{array}{lr}
&\text{\uline{Points}}\\
\text{M. . CONTOU - CARRÈRE}&=5\\
\text{J. \unsure{DEL|DE L}\unsure{i|I}MA}&=3\\
\text{O. LABORDE}&=5\\
\text{M. NG\unsure{ui|UI}FFO. BOYOM}&=7\\ \cline{2-2}
\end{array}$}
\[
\begin{aligned}
\text{Total}&=20\times25\%\\
&=\\[-.2em]
&\boxed{5\ \text{points}}
\end{aligned}
\]
\[
\left\{\begin{aligned}
\text{Valeur du point}&=670^{\mathrm f}\\
\text{Part fixe}&=7100^{\mathrm f}
\end{aligned}\right.
\]
% archive sheet number: 9
\scanpage{11}
\deleted{\unsure{Admettons}} \unsure{permettant|permettent} de les identifier à un $\overline{\mathrm m}$ $G$. On a donc

\[
\deleted{\unsure{x_2}}\left\{\begin{aligned}
&\text{ens pointé : }(G,e)\\
&\Gamma\subset G\times G\times G
\end{aligned}\right.
\]

\[
\begin{aligned}
a)\quad\forall x,y\in G,\quad &\exists\unsure{1|!}\,z\in G\ \text{tel que}\quad (x,y,z)\in\Gamma\\
&\exists1\,z\in G\quad\text{---------}\quad(z,x,y)\in\Gamma\\
&\exists1\,z\in G\quad\text{---------}\quad(x,z,y)\in\Gamma
\end{aligned}
\]

\[
\begin{aligned}
&\deleted{b)\quad\text{Soit }\unsure{x}\ \unsure{\sigma_1,\sigma_2|\sigma_1,\sigma_3}\quad\sigma_1:}
 \quad\deleted{\inserted{\sigma_1,\sigma_2}\ G\xrightarrow{\sim}G\ \text{défini par}}\\
&\deleted{(x,\sigma_1 x,e)\in\Gamma,\quad\unsure{\sigma=}}
\quad\left\{\begin{aligned}
&\deleted{(e,x,\sigma_1 x)\in\Gamma}\\
&\deleted{(\sigma_2 x,e,x)\in\Gamma}
\end{aligned}\right.
\end{aligned}
\]

$b)\quad(e,e,e)\in\Gamma$

$c)$ \unsure{Soit|Si} $\sigma_1,\sigma_2,\sigma_3:G\overset{\sim}{\text{---}}G$ défini par
\[
\begin{aligned}
(e,x,\sigma_1x)&\in\Gamma\\
(\sigma_2x,e,x)&\in\Gamma\\
(x,\sigma_3x\, e)&\in\Gamma
\end{aligned}
\]
On a \deleted{$(\sigma_3\sigma_2\sigma_1)^2=\mathrm{id}$}
\[
\sigma_2\sigma_1=\mathrm{id},\qquad\sigma_3\sigma_2=\mathrm{id},\qquad\sigma_1\sigma_3=\mathrm{id}
\]
i.e. $\sigma_1=\sigma_2=\sigma_3\ (=\sigma)\quad\sigma^2=\mathrm{id}$.

Si $x,y\in G$ soit $x\cdot y\in G$ défini par $(x,y\unsure{,|;}\sigma(xy))\in\Gamma$

On a donc $\deleted{z=xy\Longleftrightarrow (x,y}$
\[
(x,y,z)\in\Gamma\Longleftrightarrow z=\sigma(x\cdot y)
\]

% printed form
% address label upside down
MR LE PROFESSEUR GROTHENDIECK A.\\
MATHEMATIQUES\hfill 5
% archive sheet number: 10
\scanpage{12}
\mspage{- 3 -}
% One diagonal stroke crosses the typescript from above the first text line to the last column of names.
\deleted{Cette mesure aurait les avantages suivants :}

\deleted{- Coût très faible.}

\deleted{- Reclassement sur place d'un certain nombre d'Agents.}

\deleted{- La transformation des postes d'Agents de service ne gênerait en rien l'entretien de la faculté qui est effectué en grande partie par une société privé, et par du personnel hors-statut, dont l'intégration sur des postes de contractuel, type C.N.R.S. est commencée et doit se terminer d'ici 3 ans.}

\deleted{Le débat s'engage aussitôt.}

\deleted{Les Professeur présents à la réunion ainsi que le personnel pensent que cette solution est la plus rationnelle et qu'elle a les plus grandes chances de résoudre le problème.}

\deleted{Aussi,il est proposé aux présents, par les Agents de service, un vote sur la motion suivante :}

\deleted{\textquotedbl{} Les Agents de service assurant les travaux de secrétariat dans les laboratoires, les Professeurs , les représentants syndicaux et les élus au Conseil d'Université, présents à la réunion du 11 Mai 1977, demandent au Président et à l'équipe de direction de tout mettre en oeuvre afin d'obtenir les transformations des postes d'Agents de service en Agent Technique de Bureau, nécessaire au reclassement de ce personnel et au bon fonctionnement des laboratoires \textquotedbl{}.}

\deleted{Le texte de la motion est voté\unsure{.} à l'unanimité des présents.}

\deleted{Il est décidé, que cette motion sera présentée au Conseil d'Université par les représentants P.T.O.S.}

\deleted{Un collectif est choisi}\deleted{t}\deleted{ parmi les présents à la réunion. Ce collectif a pour tâche de prendre tous les contacts nécessaires au niveau de la Présidence du S.E.U., du Rectorat, afin d'obtenir effectivement les transformations de postes réclamées.}

\deleted{Il est décidé que tout le\unsure{:|.} personnel de l'U.S.T.L. sera informé de la tenue et des travaux de cette réunion et une campagne de signatures sera organisée dans un proche avenir.}

\deleted{\uline{Composition du collectif} :}

\diagram{$\begin{array}{lll}
\text{\deleted{Melle ALBOUY}}&\text{\deleted{M. GROUBERT}}&\text{\deleted{M. GIRARD}}\\
\text{\deleted{Mme \quad MALLET}}&\text{\deleted{M. ZARZYCKI}}&\text{\deleted{M. TEISSONIERE}}\\
\text{\deleted{Melle GERVA}}&\text{\deleted{M. PROUST}}&
\end{array}$}
La séance est levée à 16H 30
% archive sheet number: 11
\scanpage{13}
$(G,e)$ ens. pointé, avec loi de composition $(x,y)\mathrel{\text{---}}x\cdot y$

$\sigma$ autom de $G$ \deleted{\unsure{$,e$|$e$}} avec \quad \tikz[baseline=(v.base)]{\node (v) {$\sigma^2=\mathrm{id}$}; \draw (v.center) ellipse (2.15em and .9em);}
\[
\Gamma=\{(x\unsure{.|,}y,z)\in G^3\mid z=\sigma(x\cdot y)\}
=\{(x,y,z)\in G^3\mid x\cdot y=\sigma^{-1}z\}
\]

\uline{Conditions} \quad $a)$
\[
\begin{aligned}
\forall x,y\in G\qquad &\exists\unsure{1|!}\,z\quad\text{tel que}\quad zx=\deleted{\unsure{y}}\sigma^{-1}y\\
&\exists1\,z\quad\text{tel que}\quad xz=\sigma^{-1}y
\end{aligned}
\]
i.e. $z\mathrel{\unsure{\mapsto|\vdash}}z\deleted{\unsure{y}}x$ et $z\mathrel{\vdash}xz$ sont bijectifs

\[
\begin{aligned}
b)&\quad \sigma^{-1}e=e\cdot e\\
c)&\quad\text{On a }\forall x\quad \sigma^{-1}\sigma x=ex\qquad\text{i.e. }\deleted{\unsure{e=}}ex=x
\quad\Big\}\Rightarrow\sigma e=e\\
&\quad\forall x\quad\sigma x\cdot e=\sigma^{-1}x\qquad\text{i.e }\forall x\ \unsure{\text{on a}|\cdots}\quad x\cdot e=\sigma^{-2}x\ (\uline{=x}\\
&\quad\forall x\quad x\cdot\sigma x=\sigma^{-1}e\quad(=e)
\end{aligned}
\]

\deleted{y faisant $x=1$, on trouve $e\cdot\sigma e=\sigma^{-1}e\ (= \sigma e)$}

\deleted{donc \qquad $e\cdot\sigma$}

Donc : $e$ est unité bilatère, \deleted{et $\sigma x$ est l'inverse à droite de $x$} et les translations à g et à dr. sont bijectives. Enfin, $\forall x\in G$, $\sigma x$ est \deleted{\unsure{un inv|inverse}}\inserted{inverse} à droite \inserted{de $x$}, donc (comme $\sigma^2=\mathrm{id}$)

\[
\left.\begin{aligned}
x\cdot\sigma x&=e&(\sigma x)(\sigma\sigma x)^{\unsure{\smile}}&=e\\
\deleted{(\sigma x)x}&\deleted{{}=e}&\deleted{\sigma x\cdot x}&
\end{aligned}\right\}\quad\text{c'est un inverse bilatère}
\]

\[
\begin{aligned}
&(xy)z=x(yz)\qquad ??\ /\\
&\deleted{\unsure{\sigma x\cdot x\sigma}}\\
&\deleted{xy,z,\sigma((xy)z)}\\
&(xy,z,\sigma((xy)z))\in\Gamma\\
&(x,yz,\sigma(x(yz))\in\Gamma
\end{aligned}
\]

\[
\begin{aligned}
&(x,y,\sigma(xy))\in\Gamma\\
&\deleted{(\sigma y,\sigma x.}\\
&(x,y,z)\in\Gamma\Rightarrow z=\deleted{\unsure{(xy)\sigma(x|(x\cdot y)=\sigma x}}\sigma(x\cdot y)
\Longleftrightarrow(x\cdot y)z=\unsure{\acute e|e}\\
&\hspace{10em}\Longleftrightarrow z\cdot(x\cdot y)=e\\
&\Updownarrow\quad xy=\acute e\\
&\Updownarrow\quad(xy)z=z\\
&\phantom{\Updownarrow\quad}z(xy)=z
\end{aligned}
\]
% archive sheet number: 12
\scanpage{14}
\mspage{\textcircled{A}}
\[
\begin{aligned}
&E_1=E_{\unsure{2|1}}=E_3=G\\
&\Gamma=\{(x,y,z)\mid xyz=1\}\\
&(a,b,c)\in\Gamma\qquad\text{i.e. }\ abc=1
\end{aligned}
\]

\diagram{\begin{tikzcd}[ampersand replacement=\&,row sep=large,column sep=small]
\&E_1\arrow[dl,"{f_3}"']\arrow[dr,bend left=30,"{g_2}"]\&\\
E_2\arrow[ur,bend left=30,"{g_3}"]\arrow[rr,shift left=2pt,"{f_1}"]\&\&
E_3\arrow[ul,"{f_2}"']\arrow[ll,shift left=2pt,"{g_1}"]
\end{tikzcd}}
\[
\begin{aligned}
\deleted{f_2}\qquad f_1(x)&\quad\text{défini par}\quad (a,x,f_1(x))\in\Gamma\\
f_2(x)&\quad\text{---------}\quad(f_2(x),b,x)\in\Gamma\\
f_3(x)&\quad\text{---------}\quad(x,f_3(x),e)\in\Gamma
\end{aligned}
\]
\[
\begin{aligned}
f_1(x)&=x^{-1}a^{-1}&x^{-1}a^{-1}&=y\\
f_2(x)&=x^{-1}b^{-1}&x&=a^{-1}y^{-1}\\
f_3(x)&=x^{-1}c^{-1}&\boxed{\sigma(xy)=\sigma y\,\sigma x}
\end{aligned}
\]

\[
\begin{aligned}
&\deleted{z=xy}&&\deleted{\unsure{z}}\\
&&z=xy&\Longleftrightarrow\deleted{\unsure{x^{-1}z=y}}\\
&&\deleted{\unsure{xy=}}&x=zy^{-1}\\
&&&y^{\deleted{\unsure{+}}}=x^{-1}z
\end{aligned}
\]
\[
\begin{aligned}
g_1(x)&=f_3f_2(x)=(bx)c^{-1}\\
g_2(x)&=f_1f_3(x)=(cx)a^{-1}\\
g_3(x)&=f_2f_1(x)=(ax)b^{-1}
\end{aligned}
\]
\[
\left\{\begin{aligned}
(xy)z&=1 &y^{-1}&=zx\\
&\Updownarrow &&\Updownarrow\\
\deleted{\unsure{(yz)x}}&=1 &y&=x^{-1}z^{-1}\\
(zx)y&=1\\
&\Updownarrow\\
({\unsure{x}}^{-1}z^{-1})y^{-1}&=1\Longleftrightarrow
\end{aligned}\right.
\]

\[
\begin{aligned}
g_1g_2g_3(x)&=bcax\,b^{-1}a^{-1}c^{-1}
=(\underbrace{bca}_{\substack{\|\\1}})x(\underbrace{cab}_{\substack{\|\\1}})^{-1}=x\qquad\text{i.e. }g_1g_2g_3=\mathrm{id}\\
g_1g_2g_3&=(f_3f_2f_1)^2\\
f_3f_2f_1(x)&=(axb^{-1})^{-1}c^{-1}=bx^{-1}a^{-1}c^{-1}=bx^{-1}(ca)^{-1}=bx^{-1}b\\
\sigma_2x&=bx^{-1}b\qquad\text{de }\overline{\mathrm m}\quad\sigma_3 x=cx^{-1}c,\quad\sigma_1x=ax^{-1}a\\
&\deleted{\sigma_{\unsure{1|2}}^2 x=b(b^{-1}xb^{-1})b=x}
\end{aligned}
\]

\diagram{\begin{tikzcd}[ampersand replacement=\&,column sep=small]
E_1\& E_2\arrow[r,shift left=2pt,"f_1"]\& E_3\arrow[l,shift left=2pt,"{f_1^{-1}}"]
\end{tikzcd}}
$\deleted{(x,y,z)}$

\diagram{\begin{tikzcd}[ampersand replacement=\&,column sep=small,row sep=small]
E_1\arrow[d,"{\sigma_1}"',"{\unsure{\sim|2}}"]\&E_2\arrow[dr,"f_1"']\& E_3\arrow[dl,"{f_1^{-1}}"']\\
E_1\&E_2\&E_3
\end{tikzcd}}
\[
\begin{aligned}
&&xyz=1&\Longrightarrow\\
(x,y,z)&&ax^{-1}aa^{-1}z^{-1}y^{-1}\cdot a^{-1}&=1\\
(\sigma_1 x,f_1^{-1}z,f_1y)&=(ax^{-1}a,a^{-1}z^{-1},y^{-1}a^{-1})
\end{aligned}
\]
% archive sheet number: 13
\scanpage{15}
\mspage{\textcircled{B}}
\[
\begin{aligned}
&\deleted{\unsure{xy}}\quad z=xy && x^{-1}(uz^{-1})=y\\
u&=(xy)z\Longleftrightarrow(uz^{-1}=xy)\Longleftrightarrow (y^{-1}x^{-1})u=z\\
u&=x(yz)\Longleftrightarrow(x^{-1}u=yz)\Longleftrightarrow u(\unsure{\cdot}\\
&&x^{-1}\\
&\deleted{u=}\\
&\deleted{xy}\\
&\deleted{xy}
\end{aligned}
\]
\[
\left\{\begin{aligned}
f_1(x)&=x^{-1}a^{-1}&g_1&=f_3f_2=\deleted{\unsure{bx}}x\mathrel{\unsure{\vdash|\mapsto}}(bx)c^{-1}\\
f_2(x)&=x^{-1}b^{-1}&g_2&=f_1f_3=\qquad x\mathrel{\vdash}(cx)a^{-1}\\
f_3(x)&=x^{-1}c^{-1}&g_3&=f_2f_1=\qquad x\mathrel{\vdash}(ax)b^{-1}
\end{aligned}\right.
\]

\[
g_1g_2g_3=(f_3f_2f_1)^2\qquad x\mapsto\bigl(b\bigl(\bigl(c((ax)b^{-1})\bigr)a^{-1}\bigr)\bigr)c^{-1}
\]
\[
\begin{aligned}
\sigma_2x&=f_3f_2f_1(x)=\bigl((ax)b^{-1}\bigr)^{-1}c^{-1}
=\bigl(b(x^{-1}a^{-1})\bigr)c^{-1}\\
\sigma_2^2x&=\bigl(b\bigl(\deleted{(}\bigl[c\bigl((ax)b^{-1}\bigr)\bigr]\deleted{)}a^{-1}\bigr)\bigr)c^{-1}
\end{aligned}
\]

\[
\begin{aligned}
&x\\&ax\\&axb^{-1}\\&caxb^{-1}\\&
\deleted{b}caxb^{-1}a^{-1}\\&bcaxb^{-1}\unsure{a}\\&bcaxb^{-1}a^{-1}c^{-1}
\end{aligned}
\]
\[
\left\{\begin{aligned}
\sigma(xy)&=\sigma x\cdot\sigma y\\
\tau_x^{d}\,{\unsure{\leq|\angle}}^{-1}&=\tau_{x^{-1}}^d\\
\tau_x^{g}\,{}^{-1}&=\tau_{x^{-1}}^g .
\end{aligned}\right.
\]
\[
\begin{aligned}
&u=(xy)z\\
&\deleted{\unsure{uz^{-1}=xy}}\\
&u
\end{aligned}
\]
\[
\begin{aligned}
&\deleted{\unsure{g_1^{-1}}}\\
&(bx)c^{-1}=y\\
&x=b^{-1}(yc)
\end{aligned}
\]
\[
\begin{aligned}
&\deleted{x=}\qquad g_1g_2g_3 x=x\qquad\deleted{\unsure{g_2}g_3x=}\\
&\bigl(c((ax)b^{-1})\bigr)\deleted{a^{-1}}=\bigl(b^{-1}(xc)\bigr)a\\
&\boxed{c\bigl((ax)b^{-1}\bigr)=\bigl(b^{-1}(xc)\bigr)a}
\end{aligned}
\]

\[
\begin{aligned}
&abc\\
&\underbrace{ca}b=1\\
&\deleted{c=1}\qquad \deleted{c(xc)=(xc)}\quad\deleted{(ax)a=\unsure{x+}} \qquad \deleted{c(ab^{-1})=(b^{-1}c)a}\\
&\deleted{b=a^{-1}}\qquad\deleted{a^{-1}(ax)=(xa^{-1})a}\qquad x=
\end{aligned}
\]
% archive sheet number: 14
\scanpage{16}
\[
\boxed{(b^{-1}a^{-1})\bigl((ax)b^{-1}\bigr)
=\bigl(b^{-1}(x(b^{-1}a^{-1}))\bigr)a}
\]
\[
\boxed{\left\{\begin{aligned}
(ab)c&=1\\
&\deleted{(x\quad y\quad z)}\\
(xy)z&=1
\end{aligned}\right.\quad\Longrightarrow
\bigl[(z^{-1}b^{-1})\bigl((b(y^{-1}a^{-1}))c^{-1}\bigr)\bigr](b^{-1}x^{-1})\deleted{\unsure{\big\}|b}}=1}
\]
\[
\begin{aligned}
&\deleted{\bigl(b(x^{-1}z^{-1})\bigr)c^{-1}=b(x^{-1}b)}\\
&\underbrace{z=c^{\deleted{\unsure{-1}}}}\quad\unsure{\text{alors}}\\
&\deleted{c=b^{-1}a^{-1}}\qquad(ab)(xy)=1\\
&\deleted{z=ab.}\\
&ab=\unsure{1|/}xy\qquad\deleted{(ab)(xy)=1}\\
&ax
\end{aligned}
\]
\[
\begin{aligned}
&abc&b&\mathrel{\vdash}z\\
&c\,\unsure{{}^{\prime}|=}\,b^{-1}a^{-1}&a&\mathrel{\vdash}b\\
&\downarrow\quad\|&x&\mathrel{\unsure{\mapsto|\vdash}}y^{-1}a^{-1}\\
&\bigl[\bigl(z^{-1}\bigl((y^{-1}a^{-1})(z^{-1}b^{-1})\bigr)\bigr)b\bigr](b^{-1}x^{-1})\\
\Rightarrow{}&
\bigl[\underbrace{\bigl((\deleted{\unsure{b}}\unsure{xy})\underbrace{\bigl((y^{-1}a^{-1})\underbrace{((\unsure{ab})b^{-1})}_{a}\bigr)}_{y^{-1}}\bigr)}_{x}b\bigr](b^{-1}x^{-1})
\end{aligned}
\]
\[
\deleted{\unsure{xb}}\qquad\qquad(xb)(b^{-1}x^{-1})=1
\]
\[
\begin{aligned}
(ab)c&=1\quad *&\overset{?}{\Longrightarrow}\quad xb&=(z^{-1}b^{-1})\bigl((b(y^{-1}a^{-1}))c^{-1}\bigr)\\
(xy)z&=1\quad *&&\\
&&&\Updownarrow\\
\deleted{(ub^{-1})(bc)}&\deleted{{}=uc\ ?}&(bz)(xb)&=(b(y^{-1}a^{-1}))c^{-1}
\end{aligned}
\]
\[
\begin{aligned}
(bz)(xb)&=\bigl(b\bigl((zx)\underbrace{({\unsure{a|b}}^{-1}c)}_{\unsure{b\ ?\ b}}\bigr)\bigr)c^{-1}\\
\bigl((\underbrace{bz}_{b})(\underbrace{xb}_{\substack{\|\\\deleted{\unsure{b^{-1}a^{-1}}}\\y^{-1}}})\bigr)\underset{\substack{\|\\a^{-1}}}{c}&=b(y^{-1}a^{-1})\\
(bu)c&=b((ub^{-1})(bc))
\end{aligned}
\]
\[
\begin{aligned}
&\deleted{\unsure{b}}\\
&\boxed{\bigl((bz)(xb)\bigr)c=b\bigl((zx)(bc)\bigr)}
\qquad\left\{\begin{aligned}c&=a^{-1}\\xb&=y^{-1}\quad *\\ b&=bz\quad *\end{aligned}\right.\\
&z=1,\ c=\\
&\bigl(b(xb)\bigr)c=b\bigl(x(bc)\bigr)
\end{aligned}
\]
\[
\begin{aligned}
z&=1\qquad y={\text{\textcircled{$x$}}}^{-1}\\
b&=1,\quad c={\text{\textcircled{$a$}}}^{-1}
\end{aligned}
\]
% archive sheet number: 15
\scanpage{17}
% Dark-ink text is upside down relative to the archive number. A large pencil cross covers this layer.
\[
\begin{aligned}
&\deleted{\tau_x z=xz}\\
&\deleted{\tau_x\tau_y z=x(yz)}\\
&\deleted{xx'=x'x=e}\\
&\deleted{x(x'y)=y}\\
&\deleted{\deleted{\unsure{x}}G\times G\to G}
\end{aligned}
\]

\[
\begin{aligned}
&\deleted{e\quad a}\\
&\deleted{ee=e}\\
&\deleted{ea=ae=a}\\
&\deleted{aa\ne ae=a\quad\text{donc}\quad aa=e}\\
&\deleted{e\quad a\quad b}\\
&\deleted{\overset{\unsure{\|}}{aa}\quad bb\quad ab=ba}\\
&\deleted{b\qquad\qquad\underset{\text{\textcircled{$e$}}}{\|}\qquad\underset{e}{\|}}\\
&\deleted{\underset{\substack{\|\\a(aa)}}{ab}\overset{?}{=}\underset{\substack{\|\\(aa)a}}{ba}\ -}
\end{aligned}
\]
\[
\begin{aligned}
&\deleted{\{\underset{\substack{\|\\a}}{ae},aa,ab\}=\{e,a,b\}}\\
&\deleted{\underbrace{aa=b,\ ab=e}\qquad\text{ou}\qquad aa=e,\ ab=b}\\
&\deleted{\Downarrow\qquad\qquad\qquad\text{mais }ab\ne eb\unsure{{}=b}}\\
&\deleted{\boxed{ba=e}}\\
&\deleted{bb=a}
\end{aligned}
\]
\[
\begin{aligned}
&\deleted{e,a,b,c}&\deleted{(xy)z=}\deleted{x(yz)}&&\deleted{uv=u'v'}\\
&&\deleted{x\ y\ z}&&\deleted{\overline x((xy)z)=(yz)}\\
&\deleted{e\quad a,b,c}&&&\deleted{\overline x((xy)z)z}
\end{aligned}
\]

\diagram{$\begin{array}{c|c|c|c}
&\deleted{a}&\deleted{b}&\deleted{c}\\\hline
\deleted{a}&\deleted{e}&\deleted{\unsure{c|C}}&\deleted{b}\\\hline
\deleted{b}&&&\\\hline
\deleted{c}&&&\\\hline
\end{array}$}
\diagram{$\begin{array}{c|c|c|c|}
&\deleted{a}&\deleted{b}&\deleted{c}\\\hline
\deleted{a}&\deleted{b}&\deleted{c}&\deleted{e}\\\hline
\deleted{b}&&&\\\hline
\deleted{c}&&&\\\hline
\end{array}$}
\[
\begin{aligned}
&\deleted{aa=e\quad ab=c\quad ac=b}\\
&\deleted{a}
\end{aligned}
\]
\[
\begin{aligned}
&\deleted{(xy)z=1\Longleftrightarrow z(xy)=1}\\
&\deleted{\Updownarrow}\\
&\deleted{y(zx)=1}\\
&\deleted{\unsure{\Updownarrow}}\\
&\deleted{(zx)y=1}\\
&\deleted{(xy)z}
\end{aligned}
\]
\[
\begin{aligned}
&\deleted{(ab)c=1}&&\deleted{\unsure{\sigma x}}\\
&\deleted{(ax)y=1}&\deleted{y=\sigma(ax)}&\deleted{f_1(x)=\sigma(ax)}\\
&\deleted{(yb)x\unsure{'}=1}&\deleted{yb=\sigma x}&\deleted{f_2(x)=(\sigma x)b^{\unsure{-1|-'}}}\\
&\deleted{(x,y)c=1}&\deleted{xy=\sigma c}
\end{aligned}
\]

% Pencil entries below follow their own reading direction.
\[
(aa)a
\]
\[
\begin{aligned}
&xy=z\qquad y^{-1}=u\qquad x=zu\\
&(xy)(y^{-1}x^{-1})&a_p&=a_q a_{\unsure{r+|r|t}}\qquad\unsure{\text{prouver}}\\
&\hspace{3em}\|&&\deleted{\unsure{a_q a}}\\
&\bigl((xy)y^{-1}\bigr)x^{-1}&&\deleted{\unsure{a^p}}\\
&&&\deleted{a_p=\unsure{a_q a_{r^{-1}}}}
\end{aligned}
\]
\[
z\bigl(u(u^{-1}z^{-1})\bigr)=\bigl((zu)u^{-1}\bigr)z^{-1}
\]
% archive sheet number: 16
